\chapter{持久化、一致性与清除}

\section{身份层级}

持久化层至少区分主体标识、对象标识、修订标识、对象纪元、主体权威提交序号、运行时修订、服务视图集合、服务代次以及请求和事件的幂等标识。

对象标识表示稳定语义对象；修订标识表示不可变内容版本；对象纪元表示当前对象头、生命周期和其他可变对象状态的并发版本；权威提交序号为同一主体提供持久提交顺序。

\section{权威写入}

权威写入遵循统一的事务语义：验证请求和权限；检查幂等身份；核验预期对象纪元和当前修订；重新验证来源、访问控制和生命周期；写入不可变修订和关系；更新对象头、生命周期和纪元；记录权威变更及幂等结果；最后提交事务。

模型调用、网络搜索和其他不可控外部操作在权威写事务之外完成。它们的结果以带来源和生产者信息的提案进入事务。

同一个幂等请求重复到达时返回同一语义结果。相同幂等标识携带不同规范化内容时返回冲突，不把两次请求视为两个独立操作。

\section{投影变化与恢复}

权威提交记录需要失效或重建的服务投影类别和处理水位。投影构建器允许重试、崩溃恢复和全量重建；任何旧投影都不能覆盖较新的权威状态。

依赖其他认知对象的日记、认知图式、综合记忆、自我叙事和其他派生认知保存依赖引用。来源修订、撤回、抑制、清除或权限变化时，依赖对象按领域合同重新核验、标记冲突或进入失效状态。

\section{对象存储}

大型材料和二进制内容写入不可变对象库。写入流程先生成临时对象并校验摘要，再发布稳定对象句柄，最后在权威事务中提交引用。崩溃可以留下未引用对象并由垃圾回收清理，已提交权威引用不得指向不存在的对象。

\section{清除}

清除（Purge）由管理权威执行。清除流程先隔离普通读取并取消或标记依赖提案，然后按策略处理权威内容，再失效服务投影、缓存和受管理上下文，最后清理对象库中的相关材料并记录不可逆的清除回执。

清除完成后不得保留能够恢复被清除内容的向量、摘要或派生材料。必要的审计记录只保存操作身份、时间、范围和结果等不含被清除语义内容的信息。
